\section{来源层级与读者使用方式}

本报告把来源分为三层。第一层是公司正式披露、交易所文件与政府部门资料，用于历史财务、股本、公告时点和许可等事实。第二层是完整、可追溯的券商报告与专业商品资料，用于外部预期、价格背景和同业交叉核验。第三层是新闻、行业转述或无法取得原始底稿的线索；它们不进入基准 EPS 或目标价。

\begin{exhibitbox}[表 A-1：关键结论的证据权限]
\begin{tabularx}{\exhibitboxwidth}{L{3.0cm}X X}
\toprule
结论类型 & 可采用的证据 & 明确不采用的替代物 \\
\midrule
历史收入、利润、现金流与股本 & 公司年报、季报、业绩预告及权益分派公告 & 行情网站的财务聚合值 \\
锂盐价格与行业方向 & 具日期和单位的专业报价、官方统计与同行披露 & 用现货价格直接替代雅化实现售价 \\
客户、资源、产能与订单 & 公司或客户的正式披露、可复核合同/许可 & 媒体转述、设计产能、战略合作名称 \\
外部目标价与预测 & 完整券商 PDF 中的报告日期、方法和预测表 & 摘要、搜索片段或无日期转载 \\
自有目标价 & 本报告结构化情景模型与估值审计 & 单一券商目标、现价或 H1 利润简单年化 \\
\bottomrule
\end{tabularx}
\end{exhibitbox}

\sourcenote{公司《2025 年年度报告》《2026 年第一季度报告》《2026 年半年度业绩预告》；完整券商报告；来源登记册。}

\section{尚未披露事项的处理}

正式 H1 尚未提供分部收入、分部现金流、锂盐实际销量/实现价格/单位成本以及客户项目回款。基准情景因此不将这些项目设为已确认的超额收益来源；模型以历史已审计的利润池、审慎的价格与成本区间、以及明确的概率折扣处理。若正式 H1 的事实与此假设相反，目标价与行动将按第九章规则重估。
